\section{AMC 10 10B 2014}
        \begin{enumerate}
        	\item (\textit{AMC 10 10B 2014}) Leah has $13$ coins, all of  which are pennies and nickels. If she had one more nickel than she has now, then she would have the same number of pennies as nickels. In cents, how much are Leah's coins worth?


 $\textbf {(A) } 33 \qquad \textbf {(B) } 35 \qquad \textbf {(C) } 37 \qquad \textbf {(D) } 39 \qquad \textbf {(E) } 41$


 
 \textbf{Answer: }C



\textbf{Solution}

If Leah has $1$ more nickel, she has $14$ total coins. Because she has the same number of nickels and pennies, she has $7$ nickels and $7$ pennies. This is after the nickel has been added, so we must subtract $1$ nickel to get $6$ nickels and $7$ pennies. Therefore, Leah has $6\cdot5+7=\boxed{37 (\textbf{C})}$ cents.

 


	\item (\textit{AMC 10 10B 2014}) What is $\frac{2^3 + 2^3}{2^{-3} + 2^{-3}}$?


 $\textbf {(A) } 16 \qquad \textbf {(B) } 24 \qquad \textbf {(C) } 32 \qquad \textbf {(D) } 48 \qquad \textbf {(E) } 64$


 
 \textbf{Answer: }E



\textbf{Solution 1}

We can synchronously multiply ${2^3}$ to the expresions both above and below the fraction bar. Thus, 
 \[\frac{2^3+2^3}{2^{-3}+2^{-3}}\\=\frac{2^6+2^6}{1+1}\\={2^6}.\]Hence, the fraction equals to $\boxed{{\textbf{(E) }64}}$.

 



\textbf{Solution 2}

We have 
 \[\frac{2^3+2^3}{2^{-3}+2^{-3}} = \frac{8 + 8}{\frac{1}{8} + \frac{1}{8}} = \frac{16}{\frac{1}{4}} = 16 \cdot 4 = 64,\] so our answer is $\boxed{{(\textbf{E) }64}}$.

 


	\item (\textit{AMC 10 10B 2014}) Peter drove the first third of his trip on a gravel road, the next $20$ miles on pavement, and the remaining one-fifth on a dirt road. In miles how long was Peter's trip?


 $\textbf {(A) } 30 \qquad \textbf {(B) } \frac{400}{11} \qquad \textbf {(C) } \frac{75}{2} \qquad \textbf {(D) } 40 \qquad \textbf {(E) } \frac{300}{7}$


 
 \textbf{Answer: }E



\textbf{Solution 1}

Let the total distance be $x$. We have $\dfrac{x}{3} + 20 + \dfrac{x}{5} = x$, or $\dfrac{8x}{15} + 20 = x$. Subtracting $\dfrac{8x}{15}$ from both sides gives us $20 = \dfrac{7x}{15}$. Multiplying by $\dfrac{15}{7}$ gives us $x = \boxed{{\textbf{(E) }\dfrac{300}{7}}}$

 



\textbf{Solution 2}

The first third of his distance added to the last one-fifth of his distance equals $\frac{8}{15}$. Therefore, $\frac{7}{15}$ of his distance is $20$. Let $x$ be his total distance, and solve for $x$. Therefore, $x$ is equal to $\frac{300}{7}$, or $E$.

 


	\item (\textit{AMC 10 10B 2014}) Susie pays for $4$ muffins and $3$ bananas. Calvin spends twice as much paying for $2$ muffins and $16$ bananas. A muffin is how many times as expensive as a banana?


 $\textbf {(A) } \frac{3}{2} \qquad \textbf {(B) } \frac{5}{3} \qquad \textbf {(C) } \frac{7}{4} \qquad \textbf {(D) } 2 \qquad \textbf {(E) } \frac{13}{4}$


 
 \textbf{Answer: }B



\textbf{Solution}

Let $m$ be the cost of a muffin and $b$ be the cost of a banana. From the given information, 
 \[2m+16b=2(4m+3b)=8m+6b\Rightarrow 10b=6m\Rightarrow m=\frac{10}{6}b=\boxed{\frac{5}{3}\rightarrow \text{(B)}}\].

 


	\item (\textit{AMC 10 10B 2014}) Camden constructs a square window using $8$ equal-size panes of glass, as shown. The ratio of the height to width for each pane is $5 : 2$, and the borders around and between the panes are $2$ inches wide. In inches, what is the side length of the square window?

                    \begin{figure}[H]
                        \centering
                        
\includegraphics[width=0.5\textwidth]{Problem/AMC_10_10B_2014_5_Problem_Figure_1.png}
                    \end{figure}
                    
$\textbf{(A)}\ 26\qquad\textbf{(B)}\ 28\qquad\textbf{(C)}\ 30\qquad\textbf{(D)}\ 32\qquad\textbf{(E)}\ 34$


 
 \textbf{Answer: }A



\textbf{Solution}

We note that the total length must be the same as the total height, as it is given in the problem. Calling the width of each small rectangle $2x$, and the height $5x$, we can see that the length is composed of $4$ widths and $5$ bars of length $2$. This is equal to two heights of the small rectangles as well as $3$ bars of $2$. Thus, $4(2x) + 5(2) = 2(5x) + 3(2)$. We quickly find that $x = 2$. The total side length is  $4(4) + 5(2) = 2(10) + 3(2) = 26$, or $\boxed{\textbf{(A)}}$.

 


	\item (\textit{AMC 10 10B 2014}) Orvin went to the store with just enough money to buy $30$ balloons. When he arrived, he discovered that the store had a special sale on balloons: buy $1$ balloon at the regular price and get a second at $\frac{1}{3}$ off the regular price. What is the greatest number of balloons Orvin could buy?


 $\textbf {(A) } 33 \qquad \textbf {(B) } 34 \qquad \textbf {(C) } 36 \qquad \textbf {(D) } 38 \qquad \textbf {(E) } 39$


 
 \textbf{Answer: }C



\textbf{Solution 1}

Since he pays $\dfrac{2}{3}$ the price for every second balloon, the price for two balloons is $\dfrac{5}{3}$. Thus, if he had enough money to buy $30$ balloons before, he now has enough to buy $30 \cdot \dfrac{2}{\dfrac{5}{3}} = 30 \cdot \dfrac{6}{5} = \fbox{(C) 36}$.

 



\textbf{Solution 2}

Suppose each balloon costs $3$ dollars. Therefore, Orvin brought $90$ dollars. Since every second balloon costs $\frac{2}{3}\cdot3=2$ dollars, Orvin gets $2$ balloons for $5$ dollars. Therefore, Orvin gets $\dfrac{2\cdot90}{5}=\fbox{(C) 36}.$

 


	\item (\textit{AMC 10 10B 2014}) Suppose $A>B>0$ and $A$ is $x\%$ greater than $B$. What is $x$?


 $\textbf {(A) } 100\left(\frac{A-B}{B}\right) \qquad \textbf {(B) } 100\left(\frac{A+B}{B}\right) \qquad \textbf {(C) } 100\left(\frac{A+B}{A}\right)\qquad \textbf {(D) } 100\left(\frac{A-B}{A}\right) \qquad \textbf {(E) } 100\left(\frac{A}{B}\right)$


 
 \textbf{Answer: }A



\textbf{Solution}

We have that A is $x\%$ greater than B, so $A=\frac{100+x}{100}(B)$. We solve for $x$. We get 

 $\frac{A}{B}=\frac{100+x}{100}$

 $100\frac{A}{B}=100+x$

 $100\left(\frac{A}{B}-1\right)=x$

 $100\left(\frac{A-B}{B}\right)=x$.  $\boxed{\left(\textbf{A}\right)}$

 



\textbf{Solution 2}

The question is basically asking the percentage increase from $B$ to $A$. We know the formula for percentage increase is $\frac{\text{New-Original}}{\text{Original}}$. We know the new is $A$ and the original is $B$. We also must multiple by $100$ to get $x$ out of it's fractional/decimal form. Therefore, the answer is $100\left(\frac{A-B}{B}\right)$ or $\boxed{\text{A}}$.

 



\textbf{Solution 3 (Answer Choices)}

Without loss of generality, let $A = 125$ and $B = 100,$ forcing $x$ to be $25$. Plugging our values for $A$ and $B$ into these answer choices, we find that only $\boxed{\textbf{(A)}}$ returns $25$.

 


	\item (\textit{AMC 10 10B 2014}) A truck travels $\frac{b}{6}$ feet every $t$ seconds. There are $3$ feet in a yard. How many yards does the truck travel in $3$ minutes?


 $\textbf {(A) } \frac{b}{1080t} \qquad \textbf {(B) } \frac{30t}{b} \qquad \textbf {(C) } \frac{30b}{t}\qquad \textbf {(D) } \frac{10t}{b} \qquad \textbf {(E) } \frac{10b}{t}$


 
 \textbf{Answer: }E



\textbf{Solution}

Converting feet to yards and minutes to second, we see that the truck travels $\dfrac{b}{18}$ yards every $t$ seconds for $180$ seconds. We see that he does $\dfrac{180}{t}$ cycles of $\dfrac{b}{18}$ yards. Multiplying, we get $\dfrac{180b}{18t}$, or $\dfrac{10b}{t}$, or $\boxed{\textbf{(E)}}$.

 



\textbf{Solution 2}

We set a proportion by letting the $x$ being the number of feet the truck travels in $3$ minutes.

 \begin{align*}\frac{\frac{b}{6}}{t} &= \frac{x}{180} \\\frac{b}{6t} &= \frac{x}{180} \\\frac{180b}{6t} &= x \\x &= \frac{30b}{t}\end{align*}


 Remember $x$ is the number of feet the truck travels, so we divide by $3$ to convert to yards. 

 $\frac{x}{3}=\frac{10b}{t}$, which corresponds to $\boxed{\text{(E)}}$

 


	\item (\textit{AMC 10 10B 2014}) For real numbers $w$ and $z$, 

 \[\frac{\frac{1}{w} + \frac{1}{z}}{\frac{1}{w} - \frac{1}{z}} = 2014.\]
What is $\frac{w+z}{w-z}$?


 $\textbf{(A) } -2014 \qquad\textbf{(B) } \frac{-1}{2014} \qquad\textbf{(C) } \frac{1}{2014} \qquad\textbf{(D) } 1 \qquad\textbf{(E) } 2014$


 
 \textbf{Answer: }A



\textbf{Solution}

Multiply the numerator and denominator of the LHS (left hand side) by $wz$ to get $\frac{z+w}{z-w}=2014$. Then since $z+w=w+z$ and $w-z=-(z-w)$, $\frac{w+z}{w-z}=-\frac{z+w}{z-w}=-2014$, or choice $\boxed{A}$.

 



\textbf{Solution 2}

Basic algebra at the end of the day, so simplify the numerator and the denominator. The numerator simplifies out to $\frac{w+z}{wz}$ and the denominator simplifies out to $\frac{z-w}{wz}$. 

 This results in $\cfrac{\frac{w+z}{zw}}{\frac{z-w}{zw}} = 2014$. 

 Division results in the elimination of $zw$, so we get $\frac{w+z}{z-w} = 2014$.

 $z-w$ is just $-(w-z)$ so the equation above is $-(\frac{w+z}{w-z} = 2014$.

 Solving this results in $\frac{w+z}{w-z} = \boxed{\textbf{(A)}\ -2014}$.

 



\textbf{Solution 3}

Muliply both sides by $\left(\frac{1}{w}-\frac{1}{z}\right)$ to get $\frac{1}{w}+\frac{1}{z}=2014\left(\frac{1}{w}-\frac{1}{z}\right)$. Then, add $2014\cdot\frac{1}{z}$ to both sides and subtract $\frac{1}{w}$ from both sides to get $2015\cdot\frac{1}{z}=2013\cdot\frac{1}{w}$. Then, we can plug in the most simple values for z and w ($2015$ and $2013$, respectively), and find $\frac{2013+2015}{2013-2015}=\frac{2(2014)}{-2}=-2014$, or answer choice $\boxed{A}$.

 



\textbf{Solution 4}

Let $a = \frac{1}{w}$ and $b = \frac{1}{z}$. To find values for a and b, we can try $a+b = 2014$ and $a-b=1$. However, that leaves us with a fractional solution, so scaling it by 2, we get $a+b = 4028$ and $a-b=2$. Solving by adding the equations together, we get $b = 2015$ and $a = 2013$. Now, substituting back in, we get $w = \frac{1}{2015}$ and $z = \frac{1}{2013}$. Now, putting this into the desired equation with $n = 2015 \cdot 2013$ (since it will cancel out), we get $\frac{\frac{2013+2015}{n}}{\frac{2013-2015}{n}}$. Dividing, we get $\frac{4028}{-2} = \boxed{\textbf{(A)}\ -2014}$.

 



\textbf{Solution 5}

Set \( w = 2 \) and \( z = 1 \).

 Substitute the new values into the first equation 

 $\frac{1}{2} + 1 = \frac{3}{2}$,

 $\frac{1}{2} - 1 = -\frac{1}{2}$,

 $\frac{\frac{3}{2} / \frac{-1}{2}} = -3$

 Substitute in the second equation with new values of \( w \) and \( z \):

 $\frac{2 + 1}{2 - 1} = 3.$

 Answers of each equation (where X is the quotient): $x$ and $-x$ 

 Therefore, the answers to the equations are the negatives of each other. Thus the answer is $\boxed{(A)}$

 


	\item (\textit{AMC 10 10B 2014}) In the addition shown below $A, B, C,$ and $D$ are distinct digits. How many different values are possible for $D$?


 
 \[\begin{array}{lr}&ABBCB\\ +& BCADA\\ \hline & DBDDD\end{array}\]


 $\textbf{(A) }2\qquad\textbf{(B) }4\qquad\textbf{(C) }7\qquad\textbf{(D) }8\qquad\textbf{(E) }9$


 
 \textbf{Answer: }C



\textbf{Solution 1}

Note from the addition of the last digits that $A+B=D\text{ or }A+B=D+10$. From the addition of the frontmost digits, $A+B$ cannot have a carry, since the answer is still a five-digit number. Also $A + B$ cant have a carry since then for the second column, $C + 1 + D$ cant equal $D$. Therefore $A+B=D$. 

 Using the second or fourth column, this then implies that $C=0$, so that $B+C=B$ and $C+D=D$. Note that all of the remaining equalities are now satisfied: $A+B=D, B+C=B,$ and $B+A=D$. Since the digits must be distinct, the smallest possible value of $D$ is $1+2=3$, and the largest possible value is $9$. Thus we have that $3\le D\le9$, so the number of possible values is $\boxed{\textbf{(C) }7}$

 



\textbf{Solution 2 (Similar to 1, Less Math Jargon)}

Notice that $A$ and $B$ cannot equal $0$. 

 First, we can easily find that $C = 0$. Look at the middle column of the addition problem, which reads $B + C = B$. The only digit you can add to $B$ without changing its value is zero, so $C = 0$.

 Next, we can deduce that $A+B$ cannot be a two-digit number because it would create a carry. Look at the column immediately to the left of the middle one: $C + D = D$. Since $C = 0$, this is just $0 + D = D$. If the column to its right ($A + B$) generated a carry of $1$, this column would become $1 + 0 + D = D$, which is impossible ($1 + D \neq D$). Therefore, $A + B$ cannot carry over; it must be a single digit, meaning $A + B = D$.

 Since all digits must be distinct and $C = 0$, the smallest possible non-zero digits we can choose for $A$ and $B$ are $1$ and $2$. This means the minimum value of $A+B = D$ is $1+2=3$. Since $D$ must remain a single digit, the maximum value of $D$ is $9$. 

 Thus, there are $9 - 3 + 1$ values of $D$, which implies $\boxed{\textbf{(C) }7}$.

 


	\item (\textit{AMC 10 10B 2014}) For the consumer, a single discount of $n\%$ is more advantageous than any of the following discounts:


 (1) two successive $15\%$  discounts


 (2) three successive $10\%$ discounts


 (3) a $25\%$ discount followed by a $5\%$ discount


 What is the smallest possible positive integer value of $n$?


 $\textbf{(A)}\ \ 27\qquad\textbf{(B)}\ 28\qquad\textbf{(C)}\ 29\qquad\textbf{(D)}\ 31\qquad\textbf{(E)}\ 33$


 
 \textbf{Answer: }C



\textbf{Solution 1}

Let the original price be $x$. Then, for option $1$, the discounted price is $(1-.15)(1-.15)x = .7225x$. For option $2$, the discounted price is $(1-.1)(1-.1)(1-.1)x = .729x$. Finally, for option $3$, the discounted price is $(1-.25)(1-.05) = .7125x$. Therefore, $n$ must be greater than $\max(x - .7225x, x-.729x, x-.7125x)$. It follows $n/100$ must be greater than $.2875$. We multiply this by $100$ to get the percent value, and then round up because $n$ is the smallest integer that provides a greater discount than $28.75$, leaving us with the answer of $\boxed{\textbf{(C) } 29}$

 



\textbf{Solution 2 (a bit easier)}

Assume the original price was $100$ dollars. Thus, after a discount of $n\%$, the price will be $100-n$ dollars. Using basic calculations, find out the value of the other discounts. This leaves us with the prices: $100-n$, $\frac{289}{4}$, $\frac{9^3}{10}$, and $\frac{15\cdot19}{4}$. Simplify these to get $100-n$, $72$ (rounded down), $72.9$, and $71$ (rounded down). To have the greatest discount, we need the least price, which is $71$ dollars. Now we get the original fractional value of this, and the discount of this is $100-\frac{285}{4}$, and we round this down to $28$. Now, it's pretty easy. The integer value greater than this is $\boxed{\textbf{(C) } 29}$

 



\textbf{Solution 3 (Straightforward)}

Assume WLOG that the original price was $\$100$. Then, option 1 would cost $100 \cdot \frac{17}{20} \cdot \frac{17}{20} = \$ 72.25$. Option 2 would cost $100 \cdot \frac{9}{10} \cdot \frac{9}{10} \cdot \frac{9}{10} = \$72.90$. Option 3 would cost $100 \cdot \frac{3}{4} \cdot \frac{19}{20} = \$71.25$. Therefore, the largest integer smaller than all of these three is $71$, so $100-71= \boxed{\textbf{(C) } 29}$.

 


	\item (\textit{AMC 10 10B 2014}) The largest divisor of $2,014,000,000$ is itself. What is its fifth largest divisor?


 $\textbf{(A)}\ \ 125,875,000\qquad\textbf{(B)}\ 201,400,000\qquad\textbf{(C)}\ 251,750,000\qquad\textbf{(D)}\ 402,800,000\qquad\textbf{(E)}\ 503,500,000$


 
 \textbf{Answer: }C



\textbf{Solution}

Note that $2,014,000,000$ is divisible by $1,\ 2,\ 4,\ 5,\ 8$. Then, the fifth largest factor would come from divisibility by $8$, or $251,750,000$, or $\boxed{\textbf{(C)}}$.

 Alternative method to dividing: notice that $2,014,000,000$ factorizes into $2 \cdot 19 \cdot 53$ times $10^6$. Thus, the answer will have $7-3 = 4$ powers of 2, which means there are $4$ zeroes in the answer because each power of $2$ adds a zero.

 


	\item (\textit{AMC 10 10B 2014}) Six regular hexagons surround a regular hexagon of side length $1$ as shown. What is the area of $\triangle ABC$?


 
                    \begin{figure}[H]
                        \centering
                        \includegraphics[width=0.5\textwidth]{Problem/AMC_10_10B_2014_13_Problem_Figure_1.png}
                    \end{figure}
                    


 

$\textbf {(A) } 2\sqrt{3} \qquad \textbf {(B) } 3\sqrt{3} \qquad \textbf {(C) } 1+3\sqrt{2} \qquad \textbf {(D) } 2+2\sqrt{3} \qquad \textbf {(E) } 3+2\sqrt{3}$


 
 \textbf{Answer: }B



\textbf{Solution 1}

We note that the $6$ triangular sections in $\triangle{ABC}$ can be put together to form a hexagon congruent to each of the seven other hexagons. By the formula for the area of the hexagon, we get the area for each hexagon as $\dfrac{3\sqrt{3}}{2}$. The area of $\triangle{ABC}$, which is equivalent to two of these hexagons together, is $\boxed{\textbf{(B)}  3\sqrt{3}}$.

 



\textbf{Solution 2}

The measure of an interior angle in a hexagon is 120 degrees. Each side of the 6 triangles that make up the remainder of triangle ABC are isosceles with 2 side lengths of 1 and an angle of 120 degrees. Therefore, by the Law of Cosines, we calculate that the longest side of this triangle is $\sqrt{3}$, so the side length of triangle ABC is $2\sqrt{3}$. Using the equilateral triangle area formula, we figure out that the answer is $\boxed{\textbf{(B)}  3\sqrt{3}}$.(note that it may not be so nice to use trigonometry in AMC10 contest, however, it is a more efficient way of solving those geometry question. ~Kai Gao)

 



\textbf{Solution 3}

We know the area of a triangle can be found through the formula $\text{Area = inradius} \cdot \text{semiperimeter}$. 

 As the hexagon fully enveloped inside the triangle touches all $3$ sides, we can visualize that hexagon as $6$ congruent equilateral triangles, each with side lengths $1$. Draw a circle that circumscribes the hexagon. Using the equilateral triangles, we can see that the circle has a radius of $1$. 

 Since the circle touches all 3 sides of the triangle, we can say that $1$ is the inradius of $\triangle{ABC}$. We can find the semiperimeter of $\triangle{ABC}$ by applying the $30-60-90$ rule of triangles on the $6$ congruent triangles inside $\triangle{ABC}$ to find that the perimeter is $6\sqrt{3}$. Thus, the semiperimeter is $\dfrac{6\sqrt{3}}{2} = 3\sqrt{3}$. Thus, the area of the triangle is $1 \cdot 3\sqrt{3} = \boxed{\textbf{(B)}   3\sqrt{3}}$~NSAoPS

 


	\item (\textit{AMC 10 10B 2014}) Danica drove her new car on a trip for a whole number of hours, averaging $55$ miles per hour. At the beginning of the trip, $abc$ miles was displayed on the odometer, where $abc$ is a 3-digit number with $a \ge 1$ and $a + b + c \le 7$. At the end of the trip, the odometer showed $cba$ miles. What is $a^2 + b^2 + c^2$ ?


 $\textbf {(A) } 26 \qquad \textbf {(B) } 27 \qquad \textbf {(C) } 36 \qquad \textbf {(D) } 37 \qquad \textbf {(E) }41$


 
 \textbf{Answer: }D



\textbf{Solution 1}

Let $h$ be the number of hours Danica drove. Note that $abc$ can be expressed as $100\cdot a+10\cdot b+c$. From the given information, we have $100a+10b+c+55h=100c+10b+a$. This can be simplified into $99a+55h=99c$ by subtraction, which can further be simplified into $9a+5h=9c$ by dividing both sides by $11$. Thus we must have $h\equiv0\pmod9$. However, if $h\ge 15$, then $\text{min}\{c\}\ge\frac{9+5(15)}{9}>9$, which is impossible since $c$ must be a digit. The only value of $h$ divisible by $9$ and less than or equal to $14$ is $h=9$. 

 From this information, $9a+5(9)=9c\Rightarrow a+5=c$. Combining this with the inequalities $a+b+c\le7$ and $a\ge1$, we have $a+b+a+5\le7\Rightarrow 2a+b\le2$, which implies $a=1$, so $b=0$, and $c=6$. Thus $a^2+b^2+c^2=1+0+36=\fbox{37 \textbf{(D)}}$

 



\textbf{Solution 2}

Danica drives $m$ miles, such that $m>0$ and $m$ is a multiple of 55. Therefore, $m$ must have a units digit of either $0$ or $5.$ If the units digit of $m$ is $0,$ then $a=c$ which would imply that Danica did not drive at all. Thus, $c>a.$ Therefore, $|a-c|=5,$ and because $a+b+c\leq7, c>a,$ we have $(a,c)=(1,6).$ Finally, $b$ then must be $0$ due to $a+b+c\leq 7,$ and $a^2+b^2+c^2=1^2+0^2+6^2=\fbox{\textbf{(D) }37}$

 



\textbf{Solution 3}

We can set up an algebraic equation for this problem.

 From what's given, we have that $100c+10b+a=55x+100a+10b+c$

 This simplifies to be $0=55x+99a-99c\implies -55x=99a-99c$

 Factoring, we get that $-55x=99(a-c)\implies x=-\frac{9(a-c)}{5}$

 Hence, notice that we want $a-c=-5$ so that $x=9$

 The only pair that works for this problem that satisfies the original requirements is $(1,6)$

 Hence, $a=1, b=0, c=6$

 Checking, we have that $106+55(9)=601\implies 601=601$

 Hence, the answer is $1^2+0^2+6^2=37\implies\boxed{D}$

 


	\item (\textit{AMC 10 10B 2014}) In rectangle $ABCD$, $DC=2 \cdot CB$ and points $E$ and $F$ lie on $\overline{AB}$ so that $\overline{ED}$ and $\overline{FD}$ trisect $\angle ADC$ as shown. What is the ratio of the area of $\triangle DEF$ to the area of rectangle $ABCD$?


 
                    \begin{figure}[H]
                        \centering
                        \includegraphics[width=0.5\textwidth]{Problem/AMC_10_10B_2014_15_Problem_Figure_1.png}
                    \end{figure}
                    


 $\textbf{(A)}\ \ \frac{\sqrt{3}}{6}\qquad\textbf{(B)}\ \frac{\sqrt{6}}{8}\qquad\textbf{(C)}\ \frac{3\sqrt{3}}{16}\qquad\textbf{(D)}\ \frac{1}{3}\qquad\textbf{(E)}\ \frac{\sqrt{2}}{4}$


 
 \textbf{Answer: }A



\textbf{Solution 1}

Let the length of $AD$ be $x$, so that the length of $AB$ is $2x$ and $\text{[}ABCD\text{]}=2x^2$. 

 Because $ABCD$ is a rectangle, $\angle ADC=90^{\circ}$, and so $\angle ADE=\angle EDF=\angle FDC=30^{\circ}$. Thus $\triangle DAE$ is a $30-60-90$ right triangle; this implies that $\angle DEF=180^{\circ}-60^{\circ}=120^{\circ}$, so $\angle EFD=180^{\circ}-(120^{\circ}+30^{\circ})=30^{\circ}$. Now drop the altitude from $E$ of $\triangle DEF$, forming two $30-60-90$ triangles. 

 Because the length of $AD$ is $x$, from the properties of a $30-60-90$ triangle the length of $AE$ is $\frac{x\sqrt{3}}{3}$ and the length of $DE$ is thus $\frac{2x\sqrt{3}}{3}$. Thus the altitude of $\triangle DEF$ is $\frac{x\sqrt{3}}{3}$, and its base is $2x$, so its area is $\frac{1}{2}(2x)\left(\frac{x\sqrt{3}}{3}\right)=\frac{x^2\sqrt{3}}{3}$.

 To finish, $\frac{\text{[}\triangle DEF\text{]}}{\text{[}ABCD\text{]}}=\frac{\frac{x^2\sqrt{3}}{3}}{2x^2}=\boxed{\textbf{(A) }\frac{\sqrt{3}}{6}}$.

 



\textbf{Solution 2}

WLOG, let $AD = 1$ and $DC = 2$. Furthermore, drop an an altitude from $F$ to $CD$, which meets $CD$ at $X$. Since $\angle ADC$ is right and has been trisected, it follows that $\triangle ADE$ and $\triangle DXF$ are both $30-60-90$ triangles. Therefore, $AE = \frac{\sqrt{3}}{3}$, and $DX = AF = \sqrt{3}$. Hence, it follows that $EF = \sqrt{3} - \frac{\sqrt{3}}{3}= \frac{2\sqrt{3}}{3}$. Since $\angle ADE$ is right, the height and base of $\triangle DEF$ are $1$ and $\frac{2\sqrt{3}}{3}$, respectively. Thus, the area of $\triangle DEF$ is $\frac{\sqrt{3}}{3}$, and the area of rectengle $ABCD$ is $2$, so the ratio beween the area of $\triangle DEF$ and $ABCD$ is $\boxed{\textbf{(A) }\frac{\sqrt{3}}{6}}$. Note that we are able to assume that $AD=1$ and $DC = 2$ because we were asked to find the ratio between two areas.

 


	\item (\textit{AMC 10 10B 2014}) Four fair six-sided dice are rolled. What is the probability that at least three of the four dice show the same value?


 $\textbf {(A) } \frac{1}{36} \qquad \textbf {(B) } \frac{7}{72} \qquad \textbf {(C) } \frac{1}{9}\qquad \textbf {(D) } \frac{5}{36} \qquad \textbf {(E) } \frac{1}{6}$


 
 \textbf{Answer: }B



\textbf{Solution 1}

We split this problem into $2$ cases.

 First, we calculate the probability that all four are the same. After the first dice, all the numbers must be equal to that roll, giving a probability of $1 \cdot \dfrac{1}{6} \cdot \dfrac{1}{6} \cdot \dfrac{1}{6} = \dfrac{1}{216}$.

 Second, we calculate the probability that three are the same and one is different. After the first dice, the next two must be equal and the third different. There are $4$ orders to roll the different dice, giving $4 \cdot 1 \cdot \dfrac{1}{6} \cdot \dfrac{1}{6} \cdot \dfrac{5}{6} = \dfrac{5}{54}$.

 Adding these up, we get $\dfrac{7}{72}$, or $\boxed{\textbf{(B)}}$.

 


	\item (\textit{AMC 10 10B 2014}) What is the greatest power of $2$ that is a factor of $10^{1002} - 4^{501}$?


 $\textbf{(A) } 2^{1002} \qquad\textbf{(B) } 2^{1003} \qquad\textbf{(C) } 2^{1004} \qquad\textbf{(D) } 2^{1005} \qquad\textbf{(E) }2^{1006}$


 
 \textbf{Answer: }D



\textbf{Solution 1}

We begin by factoring the $2^{1002}$ out. This leaves us with $5^{1002} - 1$. 

 We factor the difference of squares, leaving us with $(5^{501} - 1)(5^{501} + 1)$. We note that all even powers of $5$ more than two end in ...$625$. Also, all odd powers of five more than $2$ end in  ...$125$. Thus, $(5^{501} + 1)$ would end in ...$126$ and thus would contribute one power of two to the answer, but not more. 

 We can continue to factor $(5^{501} - 1)$ as a difference of cubes, leaving us with $(5^{167} - 1)$ times an odd number (Notice that the other number is $5^{334} + 5^{167} + 1$. The powers of $5$ end in $5$, so the two powers of $5$ will end with $0$. Adding $1$ will make it end in $1$. Thus, this is an odd number). $(5^{167} - 1)$ ends in ...$124$, contributing two powers of two to the final result.

 Or we can see that $(5^{501} - 1)$ ends in $124$, and is divisible by $2$ only. Still that's $2$ powers of $2$.

 Adding these extra $3$ powers of two to the original $1002$ factored out, we obtain the final answer of $\textbf{(D) } 2^{1005}$.

 



\textbf{Solution 2}

First, we can write the expression in a more primitive form which will allow us to start factoring.
 \[10^{1002} - 4^{501} = 2^{1002} \cdot 5^{1002} - 2^{1002}\]Now, we can factor out $2^{1002}$. This leaves us with $5^{1002} - 1$. Call this number $N$. Thus, our final answer will be $2^{1002+k}$, where $k$ is the largest power of $2$ that divides $N$. Now we can consider $N \pmod{16}$, since $k \le 4$ by the answer choices. 

 Note that
 \begin{align*} 5^1 &\equiv 5 \pmod{16} \\ 5^2 &\equiv 9 \pmod{16} \\ 5^3 &\equiv 13 \pmod{16} \\ 5^4 &\equiv 1 \pmod{16} \\ 5^5 &\equiv 5 \pmod{16} \\ &\: \: \qquad \vdots \end{align*}
The powers of $5$ cycle in $\mod{16}$ with a period of $4$. Thus, 
 \[5^{1002} \equiv 5^2 \equiv 9 \pmod{16} \implies 5^{1002} - 1 \equiv 8 \pmod{16}\]This means that $N$ is divisible by $8= 2^3$ but not $16 = 2^4$, so $k = 3$ and our answer is $2^{1002 + 3} =\boxed{\textbf{(D)}\: 2^{1005}}$.

 



\textbf{Solution 3}

Convert $4^{501}=2^{1002}$. We can factor out $2^{1002}$ to get that $\nu_2(10^{1002}-2^{1002})=1002+\nu_2(5^{1002}-1)$. Using the adjusted Lifting The Exponent lemma ($\nu_2(a^n-b^n)=\nu_2(n)+\nu_2(a^2-b^2)-1$ for all even $n$ and odd $a,b$), we get that the answer is $2^{1002+\nu_2(1002)+\nu_2(24)-1}=2^{1002+1+3-1}=\boxed{\textbf{(D)}2^{1005}}$

 



\textbf{Solution 4}

Factor out $2^{1002}$ to get $2^{1002}(5^{1002} - 1)$. Since $5^{1002}-1\equiv (-3)^{1002}-1\equiv (9)^{501}-1\equiv 1^{501}-1\equiv 0\pmod{8}$, but $5^{1002}-1\equiv 11^{1002}-1\equiv 121\cdot (11^4)^{250} - 1\equiv 121 - 1\equiv 8 \pmod{16}$, $5^{1002}-1$ has 3 factors of 2. Hence $2^{1002 + 3} =\boxed{2^{1005}(D)}$ is the largest power of two which divides the given number

 



\textbf{Solution 5}

Like Solution 1, factor out $2^{1002}$ to get $(2^{1002})(5^{501}-1)(5^{501}+1)$. Using engineer's induction, we observe that for any positive integer $5^n$ (where $n$ is an odd positive integer), it appears that the least even numbers directly above and below $n$ in value must contain a maximum multiple of $4$ and a maximum multiple of $2$. Hence, the answer is $2^{1002+2+1}$ which is $\boxed{\textbf{(D)} 2^{1005}}$ .

 Proof;

 For all integers $x$ where $x=5^{n}$ where n is an odd integer, $x$ must end in $125.$ Thus, we find that $x-1$ and $x+1$ respectively end in $124$ and $126.$ 

 Case $1$ : $x-1$

 We know that this number takes the form $abcde... 124$ where $abcde...$ is an integer that ends in $124$. Because $abcde...$ is a multiple of $4$ times an even number $e$ while $124$ is $4(31)$, we find that $X-1$ must be $4e + 4 \cdot 31 = 4(e+31) = 4o$ where $o$ is an odd number

 Case $2$ : $x+1$

 We know that this number $fghijk...126$ ends. Because it is $2$ more than the number $x-1$, which is a multiple of $4$, we find $x+1 = 4o + 2$ which is an even number that is not divisible by $4$. Thus, it must have a maximum of $1$ multiple of $2$.

 This means that for any number $x$ being in the form $5^{n}$ where $n$ is an odd integer, $x-1$ must have a maximum of $2$ factors of $2$ while $x+1$ must have a maximum of $1$ factor of $2$. 

 


 



\textbf{Solution 6}

Using difference of squares, we get $\left(10^{501}-2^{501}\right) \left(10^{501}+2^{501}\right)$. Factoring a $2^{501}$ out, we get $\left(2^{501}\right) \left(5^{501}-1\right) \left(2^{501}\right) \left(5^{501}+1\right)$, and grouping like terms give $\left(2^{1002}\right) \left(5^{501}-1\right) \left(5^{501}+1\right)$. 

 
Then, you would go ahead and innocently choose $\textbf{(A) } 2^{1002}$, right? No! Note that $5^n$, where $n$ is any odd integer greater than or equal to $3$, it always ends in $125$. So, $5^{501}+1$ ends in $126$ and $5^{501}-1$ ends in $124$, so they add up to an extra three $2$'s. Therefore, the answer is actually $2^{1002+3}=\boxed{\textbf{(D) } 2^{1005}}$.

 



\textbf{Solution 7}

Note that we are trying to find the number of powers of $2$ in $(5^{1002}-1) \cdot 2^{1002}.$ Let's see how many $5^{1002}-1$ has.

 Try the first 4 powers of 5, namely $5^1$, $5^2$, $5^3$, and $5^4$. Note that when taking mod 4, all result in $1 \pmod 4$. When taking mod 8, even powers result in $1 \pmod 8$. When taking mod 16, every $4$th power (i.e. $4,8,..$) will result in $1 \pmod{16}$. 

 Because $1002$ is even but $2 \pmod 4, 5^{1002}$ will be equivalent to $1 \pmod 8$ but not $1 \pmod 16$. Hence $5^{1002} - 1$ \equiv $0 \pmod 8$, and so we have an extra power of $2^3$, hence the result is $1002+3=1005$. Hence the answer is (D).

 



\textbf{Solution 8}

We can rewrite the expression as $2^{1002} \cdot 5^{1002} - 2^{1002}$. Then factoring out a $2^{1002}$ gets us $5^{1002}-1$. Then we can factor it by using the difference of squares formula. $(5^{501}-1)(5^{501}+1)$. $(5^{501}+1)$ can be factored out further by using the difference of cubes formula, $5^{501}+1 = \left(5^{167}\right)^3 + 1^3 = \left(5^{167}+1\right)\left(5^{334} - 5^{167} + 1\right)$. We see that $\left(5^{167}\right)^3 + 1^3$ ends with 6, so dividing by 2 gets a number that ends with 3, so $5^{501} +1$ gets us 1 factor of 2. Now we deal with $\left(5^{334} - 5^{167} + 1\right)$. $5^{334}, 5^{167},$ and $1$ are all odd numbers, and the sum of 3 odd numbers is odd, so there are no powers of 2 from $\left(5^{334} - 5^{167} + 1\right)$. $5^{501}-1$ can be factored into $5^{501}-1 = (5-1)\left(5^{500} + 5^{499} + \cdots + 5^1 + 5^0\right)$. 

 $5-1=4$ gets us 2 powers of 2. $\left(5^{500} + 5^{499} + \cdots + 5^1 + 5^0\right)$ is the sum of 501 odd numbers, so that is odd, which gets us no powers of 2. Therefore, in total, $2^{1002} \cdot 5^{1002} - 2^{1002}$ has $1002+1+2$ factors of 2, which gets us the answer $\boxed{(D)}$

 


	\item (\textit{AMC 10 10B 2014}) A list of $11$ positive integers has a mean of $10$, a median of $9$, and a unique mode of $8$. What is the largest possible value of an integer in the list?


 $\textbf {(A) } 24 \qquad \textbf {(B) } 30 \qquad \textbf {(C) } 31\qquad \textbf {(D) } 33 \qquad \textbf {(E) } 35$


 
 \textbf{Answer: }E



\textbf{Solution 1}

We start off with the fact that the median is $9$, so we must have $a, b, c, d, e, 9, f, g, h, i, j$, listed in ascending order. Note that the integers do not have to be distinct. 

 Since the mode is $8$, we have to have at least $2$ occurrences of $8$ in the list. If there are $2$ occurrences of $8$ in the list, we will have $a, b, c, 8, 8, 9, f, g, h, i, j$. In this case, since $8$ is the unique mode, the rest of the integers have to be distinct. So we minimize $a,b,c,f,g,h,i$ in order to maximize $j$. If we let the list be $1,2,3,8,8,9,10,11,12,13,j$, then $j = 11 \times 10 - (1+2+3+8+8+9+10+11+12+13) = 33$. 

 Next, consider the case where there are $3$ occurrences of $8$ in the list. Now, we can have two occurrences of another integer in the list. We try $1,1,8,8,8,9,9,10,10,11,j$. Following the same process as above, we get $j = 11 \times 10 - (1+1+8+8+8+9+9+10+10+11) = 35$. As this is the highest choice in the list, we know this is our answer. Therefore, the answer is $\boxed{\textbf{(E) }35}$

 


 



\textbf{Solution 2}

Note that $x_1 + \ldots + x_{11} = 110$ let $x_6 = 9$ so $x_1 + \ldots + x_5 + x_7 + \ldots + x_{11} = 101$. To maximize the value of $x_i$ where $i$ ranges from $1$ to $11$, we let any $7$ elements be $1,2,\ldots,7$ so $x_1 + x_2 + x_3 = 57$. Now we have to let one of above $3$ values = $8$ hence $x_1 + x_2 = 49$ now let $x_1 = 35$, $x_2 = 14$ hence $\boxed{\textbf{(E) }35}$ is the answer.

 


	\item (\textit{AMC 10 10B 2014}) Two concentric circles have radii $1$ and $2$. Two points on the outer circle are chosen independently and uniformly at random. What is the probability that the chord joining the two points intersects the inner circle?


 $\textbf{(A) }\frac{1}{6}\qquad\textbf{(B) }\frac{1}{4}\qquad\textbf{(C) }\frac{2-\sqrt{2}}{2}\qquad\textbf{(D) }\frac{1}{3}\qquad\textbf{(E) }\frac{1}{2}\qquad$


 
 \textbf{Answer: }D



\textbf{Solution}

Let the center of the two circles be $O$. Now pick an arbitrary point $A$ on the boundary of the circle with radius $2$. We want to find the range of possible places for the second point, $A'$, such that $AA'$ passes through the circle of radius $1$. To do this, first draw the tangents from $A$ to the circle of radius $1$. Let the intersection points of the tangents (when extended) with circle of radius $2$ be $B$ and $C$. Let $H$ be the foot of the altitude from $O$ to $\overline{BC}$. Then we have the following diagram.

 
                    \begin{figure}[H]
                        \centering
                        \includegraphics[width=0.5\textwidth]{Solution/AMC_10_10B_2014_19_Solution_Figure_1.png}
                    \end{figure}
                    

 We want to find $\angle BOC$, as the range of desired points $A'$ is the set of points on minor arc $\overarc{BC}$. This is because $B$ and $C$ are part of the tangents, which "set the boundaries" for $A'$. Since $OH = 1$ and $OB = 2$ as shown in the diagram, $\triangle OHB$ is a $30-60-90$ triangle with $\angle BOH = 60^\circ$. Thus, $\angle BOC = 120^\circ$, and the probability $A'$ lies on the minor arc $\overarc{BC}$ is thus $\dfrac{120}{360} = \boxed{\textbf{(D)}\: \dfrac13}$.

 \textbf{Comment} See talk page

 


	\item (\textit{AMC 10 10B 2014}) For how many integers $x$ is the number $x^4 - 51x^2 + 50$ negative?


 $\textbf {(A) } 8 \qquad \textbf {(B) } 10 \qquad \textbf {(C) } 12\qquad \textbf {(D) } 14 \qquad \textbf {(E) }16$


 
 \textbf{Answer: }C



\textbf{Solution 1}

First, note that $50+1=51$, so we factor the polynomial as $(x^2-50)(x^2-1)$. 

 Since this expression is negative, one term must be negative and the other positive. We note that the first term is a lot smaller than the second, so we get $x^2-50<0<x^2-1$. 

 Solving this inequality, we find $1<x^2<50$. There are  exactly $12$ integers $x$ that satisfy this inequality, $\pm \{2,3,4,5,6,7\}$.

 Thus our answer is $\boxed{\textbf {(C) } 12}$.

 



\textbf{Solution 2}

Since the $x^4-51x^2$ part of $x^4-51x^2+50$ has to be less than $-50$ (because we want $x^4-51x^2+50$ to be negative), we have the inequality $x^4-51x^2<-50 \rightarrow x^2(x^2-51) <-50$. $x^2$ has to be positive, so $(x^2-51)$ is negative. Then we have $x^2<51$. We know that if we find a positive number that works, it's parallel negative will work. Therefore, we just have to find how many positive numbers work, then multiply that by $2$. If we try $1$, we get $1^4-51(1)^4+50 = -50+50 = 0$, and $0$ therefore doesn't work. Test two on your own, and then proceed. Since two works, all numbers above $2$ that satisfy $x^2<51$ work, that is the set {${2,3,4,5,6,7}$}. That equates to $6$ numbers. Since each numbers' negative counterparts work, $6\cdot2=\boxed{\textbf{(C) }12}$.

 



\textbf{Solution 3 (Graph)}

As with Solution $1$, note that the quartic factors to $(x^2-50)\cdot(x^2-1)$, which means that it has roots at $-5\sqrt{2}$, $-1$, $1$, and $5\sqrt{2}$. Now, because the original equation is of an even degree and has a positive leading coefficient, both ends of the graph point upwards, meaning that the graph dips below the $x$-axis between $-5\sqrt{2}$ and $-1$ as well as $1$ and $5\sqrt{2}$. $5\sqrt{2}$ is a bit more than $7$ ($1.4\cdot 5=7$) and therefore means that $-7,-6,-5,-4,-3,-2,2,3,4,5,6,7$ all give negative values.

 



\textbf{Solution 4}

Let $x^{2}=u$. Then the expression becomes $u^{2}-51u+50$ which can be factored as $\left(u-1\right)\left(u-50\right)$. Since the expression is negative, one of $\left(u-1\right)$ and $\left(u-50\right)$ need to be negative. $u-1>u-50$, so $u-1>0$ and $u-50<0$, which yields the inequality $50>u>1$. Remember, since $u=x^{2}$ where $x$ is an integer, this means that $u$ is a perfect square between $1$ and $50$. There are $6$ values of $u$ that satisfy this constraint, namely $4$, $9$, $16$, $25$, $36$, and $49$. Solving each of these values for $x$ yields $12$ values (as $x$ can be negative or positive) $\Longrightarrow \boxed{\textbf {(C) } 12}$.~JH. L

 


 



\textbf{Solution 5 (Completing the Square)}

Let $a = x^2$. We want to find when the expression is negative, so let's define the quadratic function:
 \[y = a^2 - 51a + 50\]

 We can complete the square to rewrite this function in vertex form:
 \[y = \left(a - \frac{51}{2}\right)^2 - \frac{2601}{4} + 50\]
 \[y = \left(a - \frac{51}{2}\right)^2 - \frac{2401}{4}\]

 This represents a parabola that opens upwards with its vertex located at:
 \[\left(\frac{51}{2}, -\frac{2401}{4}\right)\]

 Since the vertex lies below the axis, the function $y$ is negative ($y < 0$) between its two boundary intercepts. To find where the parabola crosses the axis, we set $y = 0$:
 \[\left(a - \frac{51}{2}\right)^2 - \frac{2401}{4} = 0\]
 \[\left(a - \frac{51}{2}\right)^2 = \frac{2401}{4}\]

 Taking the square root of both sides gives:
 \[a - \frac{51}{2} = \pm \frac{49}{2}\]
 \[a = \frac{51}{2} \pm \frac{49}{2}\]

 This gives our two boundary intercepts for $a$:

 
                \begin{itemize}
                    \item $a = \frac{51}{2} + \frac{49}{2} = \frac{100}{2} = 50$\item $a = \frac{51}{2} - \frac{49}{2} = \frac{2}{2} = 1$
                \end{itemize}


                Therefore, the function is negative when $1 < a < 50$. Substituting $a = x^2$ back in:
 \[1 < x^2 < 50\]

 The integer values for $x$ that satisfy this are $x = \pm 2, \pm 3, \pm 4, \pm 5, \pm 6, \pm 7$. Counting these values gives $6 \times 2 = 12$ possible integers $\implies \boxed{\textbf{(C)}\ 12}$.

 


	\item (\textit{AMC 10 10B 2014}) Trapezoid $ABCD$ has parallel sides $\overline{AB}$ of length $33$ and $\overline{CD}$ of length $21$. The other two sides are of lengths $10$ and $14$. The angles at $A$ and $B$ are acute. What is the length of the shorter diagonal of $ABCD$?


 $\textbf{(A) } 10\sqrt{6} \qquad\textbf{(B) } 25 \qquad\textbf{(C) } 8\sqrt{10} \qquad\textbf{(D) } 18\sqrt{2} \qquad\textbf{(E) } 26$


 
 \textbf{Answer: }B



\textbf{Solution 1}


                    \begin{figure}[H]
                        \centering
                        \includegraphics[width=0.5\textwidth]{Solution/AMC_10_10B_2014_21_Solution_Figure_1.png}
                    \end{figure}
                    

 In the diagram, $\overline{DE} \perp \overline{AB}, \overline{FC} \perp \overline{AB}$. Denote $\overline{AE} = x$ and $\overline{DE} = h$. In right triangle $AED$, we have from the Pythagorean theorem: $x^2+h^2=100$. Note that since $EF = DC$, we have $BF = 33-DC-x = 12-x$. Using the Pythagorean theorem in right triangle $BFC$, we have $(12-x)^2 + h^2 = 196$. 

 
We isolate the $h^2$ term in both equations, getting $h^2= 100-x^2$ and$h^2 = 196-(12-x)^2$.

 Setting these equal, we have $100-x^2 = 196 - 144 + 24x -x^2 \implies 24x = 48 \implies x = 2$. Now, we can determine that $h^2 = 100-4 \implies h = 4\sqrt{6}$.

 
                    \begin{figure}[H]
                        \centering
                        \includegraphics[width=0.5\textwidth]{Solution/AMC_10_10B_2014_21_Solution_Figure_2.png}
                    \end{figure}
                    

 The two diagonals are $\overline{AC}$ and $\overline{BD}$. Using the Pythagorean theorem again on $\bigtriangleup AFC$ and $\bigtriangleup BED$, we can find these lengths to be $\sqrt{96+529} = 25$ and $\sqrt{96+961} = \sqrt{1057}$. Since $\sqrt{96+529}<\sqrt{96+961}$, $25$ is the shorter length*, so the answer is $\boxed{\textbf{(B) }25}$.

 
                \begin{itemize}
                    \item Or, alternatively, one can notice that the two triangles have the same height but $\bigtriangleup AFC$ has a shorter base than $\bigtriangleup BED$.
                \end{itemize}


                



\textbf{Solution 2}


                    \begin{figure}[H]
                        \centering
                        \includegraphics[width=0.5\textwidth]{Solution/AMC_10_10B_2014_21_Solution_Figure_3.png}
                    \end{figure}
                    

 The area of $\Delta AED$ is by Heron's, $4\sqrt{9(4)(3)(2)}=24\sqrt{6}$. This makes the length of the altitude from $D$ onto $\overline{AE}$ equal to $4\sqrt{6}$. One may now proceed as in Solution $1$ to obtain an answer of $\boxed{\textbf{(B) }25}$.

 


 



\textbf{Solution 3}

Using the same way as Solution 1, obtain that $AE=2$. Extend $DC$ and drop a perpendicular from $A$ onto $DC$. Call this point $G$. We know that $GC=DC+AE=21+2=23$, and following the method from Solution 1, we also have $AG=4\sqrt{6}$. Thus the answer would be the 
 \[\sqrt{23^2+(4\sqrt{6})^2}=\sqrt{529+96}=\sqrt{625}=\boxed{\textbf{(B) } 25}\]-Reality Writes

 


	\item (\textit{AMC 10 10B 2014}) Eight semicircles line the inside of a square with side length 2 as shown. What is the radius of the circle tangent to all of these semicircles?


 $\text{(A) } \dfrac{1+\sqrt2}4 \quad \text{(B) } \dfrac{\sqrt5-1}2 \quad \text{(C) } \dfrac{\sqrt3+1}4 \quad \text{(D) } \dfrac{2\sqrt3}5 \quad \text{(E) } \dfrac{\sqrt5}3$


 
                    \begin{figure}[H]
                        \centering
                        \includegraphics[width=0.5\textwidth]{Problem/AMC_10_10B_2014_22_Problem_Figure_1.png}
                    \end{figure}
                    


 
 \textbf{Answer: }B



\textbf{Solution}

We connect the centers of the circle and one of the semicircles, then draw the perpendicular from the center of the middle circle to that side, as shown.

 
                    \begin{figure}[H]
                        \centering
                        \includegraphics[width=0.5\textwidth]{Solution/AMC_10_10B_2014_22_Solution_Figure_1.png}
                    \end{figure}
                    

 We will start by creating an equation by the Pythagorean theorem: 
 \[\sqrt{1^2 + \left(\frac12\right)^2} = \sqrt{\frac54} = \frac{\sqrt5}{2}.\]

 
Let's call $r$ as the radius of the circle that we want to find. We see that the hypotenuse of the bold right triangle is $\dfrac{1}{2}+r$, and thus $r$ is $\boxed{\textbf{(B)} \frac{\sqrt{5}-1}{2}}$.

 -- LORD\_ERTY09

 


 


	\item (\textit{AMC 10 10B 2014}) A sphere is inscribed in a truncated right circular cone as shown. The volume of the truncated cone is twice that of the sphere. What is the ratio of the radius of the bottom base of the truncated cone to the radius of the top base of the truncated cone?


 
                    \begin{figure}[H]
                        \centering
                        \includegraphics[width=0.5\textwidth]{Problem/AMC_10_10B_2014_23_Problem_Figure_1.png}
                    \end{figure}
                    


 $\text{(A) } \dfrac32 \quad \text{(B) } \dfrac{1+\sqrt5}2 \quad \text{(C) } \sqrt3 \quad \text{(D) } 2 \quad \text{(E) } \dfrac{3+\sqrt5}2$


 
 \textbf{Answer: }E



\textbf{Solution 1}

First, we draw the vertical cross-section passing through the middle of the frustum.Let the top base have a diameter of $2$ and the bottom base has a diameter of $2r$.

 
                    \begin{figure}[H]
                        \centering
                        \includegraphics[width=0.5\textwidth]{Solution/AMC_10_10B_2014_23_Solution_Figure_1.png}
                    \end{figure}
                    

 Then using the Pythagorean theorem we have:$(r+1)^2=(2s)^2+(r-1)^2$, which is equivalent to:$r^2+2r+1=4s^2+r^2-2r+1$.Subtracting $r^2-2r+1$ from both sides,$4r=4s^2$, and solving for $s$, we end up with
 \[s=\sqrt{r}.\]Next, we can find the volume of the frustum (truncated cone) and the sphere. Since we know $V_{\text{frustum}}=2V_{\text{sphere}}$, we can solve for $s$using $V_{\text{frustum}}=\frac{\pi h}{3}(R^2+r^2+Rr)$we get:
 \[V_{\text{frustum}}=\frac{\pi\cdot2\sqrt{r}}{3}(r^2+r+1)\]Using  $V_{\text{sphere}}=\dfrac{4s^{3}\pi}{3}$, we get 
 \[V_{\text{sphere}}=\dfrac{4(\sqrt{r})^{3}\pi}{3}\]so we have:
 \[\frac{\pi\cdot2\sqrt{r}}{3}(r^2+r+1)=2\cdot\dfrac{4(\sqrt{r})^{3}\pi}{3}.\]Dividing by $\frac{2\pi\sqrt{r}}{3}$, we get
 \[r^2+r+1=4r\]which is equivalent to 
 \[r^2-3r+1=0\]by the Quadratic Formula, $r=\dfrac{3\pm\sqrt{(-3)^2-4\cdot1\cdot1}}{2\cdot1}$, so
 \[r=\dfrac{3+\sqrt{5}}{2} \longrightarrow \boxed{\textbf{(E)}}\]

 



\textbf{Solution 2}

Similar to above, draw a smaller cone top with the base of the smaller circle with radius $r_1$ and height $h$. The smaller right triangle is similar to the blue highlighted one in Solution $1$. Then $\frac{r_1}{r_2-r_1}=\frac{h}{2R}$ where $R$ is the radius of the sphere. Then $h=\frac{2Rr_1}{r_2-r_1}$.

 From the Pythagorean theorem on the blue triangle in Solution $1$, we get similarly that $R=\sqrt{r_2r_1}$.

 From the volume requirements, we get that $\frac{8}{3}\pi R^3=\frac{\pi r_2^2(h+2R)-\pi r_1^2h}{3}$ which yields $8R^3=r_2^2(h+2R)-r_1^2h$.

 The small right triangle on top is similar to the big right triangle of the entire big cone. So $\frac{r_1}{r_2}=\frac{h}{h+2R}\implies h+2R=\frac{r_2h}{r_1}$. 

 Substituting yields $8R^3=\frac{h(r_2^3-r_1^3)}{r_1}$. 

 Substituting $R=\sqrt{r_2r_1}$ and $h=\frac{2Rr_1}{r_2-r_1}$ yields $8(r_2r_1)^{\frac{3}{2}}=2(r_2r_1)^{\frac{1}{2}}(r_2^2+r_2r_1+r_1^2)$ which yields $r_2^2-3r_2r_1+r_1^2=0$.

 Solving in $r_2$ yields $r_2=\frac{3r_1\pm r_1\sqrt{5}}{2}$ so $\frac{r_2}{r_1}=\frac{3+\sqrt{5}}{2}$.

 



\textbf{Solution 3 (Safe Risks)}

We can create a cross section right down the middle vertically of the truncated cone and sphere. The figure formed (excuse my horrendous editing skills) looks like below.

 
                    \begin{figure}[H]
                        \centering
                        \includegraphics[width=0.5\textwidth]{Solution/AMC_10_10B_2014_23_Solution_Figure_2.png}
                    \end{figure}
                    We can conjoin the bottom most points to form a kite. We can now use safe assumptions, a very common competition math technique that involves making estimated guesses to assume an answer. The kite has the shorter isosceles parts as length 3. We can assume the line \( AM = 3 \), and therefore line \( IE = 6 \). Then, triangle \( AIE \) is a 30-60-90 right angle triangle, and \( AE = 3\sqrt{3} \), which means that the radius of the bottom base of the truncated cone is \( 3\sqrt{3} \). We may apply the same logic to the top part of the trapezoid. Like the bottom, the length \( IM = IC = 3 \), and \( ID = \) some value that makes \( ICD \) a 30-60-90 right triangle, and therefore length \( CD = \sqrt{3} \). 

 Our ratio for the big radius to the small radius is \( 3\sqrt{3}/\sqrt{3} = 3 \). However, 3 isn't an answer :-(, so we do some educated guessing. For simplicity we say that \( \sqrt{5} = 2 \). We see that answer choice A is 3/2, B is also 3/2, C is \(\sqrt{3}\), D is 2, and E is 2.5. We see that E is closest to the ratio we got, and so our answer is $\boxed{\textbf{(E)}}$.

 



\textbf{Solution 4: Another Way of Simplifying the Equation in Solution 2}

Similar to the previous solutions, we let $R$, $r$, and $s$ be the radii of the larger base, smaller base, and the sphere, respectively. By the Pythagorean Theorem, we have $s=\sqrt{Rr}$. Let $H$, $h$ be the heights of the large and small cones, respectively, then we have $H=h+2s$ and
 \[h=2s\left(\frac{r}{R-r}\right)==\frac{2sr}{R-r}=\frac{2\left(\sqrt{Rr}\right)r}{R-r}\] Thus we have the equation 
 \[V_{large cone}-V_{small cone}=2V_{sphere}\]
 \[\frac{\pi R^2H}{3}-\frac{\pi r^2h}{3}=2\left(\frac{4\pi s^3}{3}\right)\]
 \[\pi R^2H-\pi r^2h=8\pi s^3\]
 \[R^2H-r^2h=8s^3\]Substituting the expressions for $H$, $h$, $s$ in order into the equation yields
 \[R^2\left(h+2s\right)+r^2h=8s^3\]
 \[R^2\left(\left(\frac{2sr}{R-r}\right)+2s\right)+r^2\left(\frac{2sr}{R-r}\right)=8s^3\]
 \[R^2\left(\frac{2sR}{R-r}\right)+r^2\left(\frac{2sr}{R-r}\right)=8s^3\]
 \[\frac{2sR^3-2sr^3}{R-r}=8s^3\]
 \[\frac{R^3-r^3}{R-r}=4s^2\]The numerator $R^3-r^3$ can be factored to yield
 \[\frac{\left(R-r\right)\left(R^2+Rr+r^2\right)}{R-r}=4s^2\]
 \[R^2+Rr+r^2{R-r}=4s^2\]
 \[R^2+Rr+r^2=4\left(\sqrt{Rr}\right)^2\]
 \[R^2+Rr+r^2=4Rr\]
 \[R^2-3Rr+r^2=0\]Solving the quadratic to obtain the ratio $\frac{R}{r}=\frac{3+\sqrt{5}}{2}$ $\boxed{\textbf{(E)}}$

 


 



\textbf{Solution 5 (Trapezoid Tangency Property)}

Let the frustum have bottom radius $R$, top radius $r$, and height $h$. Let the inscribed sphere have radius $a$.

 Take a vertical cross-section through the center. This shows an isosceles trapezoid with bases $2R$ and $2r$, height $h$, and an inscribed circle of radius $a$.

 Because the circle touches both bases, the frustum height equals the sphere’s diameter:
 \[h = 2a\]

 The slant height $l$ is the hypotenuse of the right triangle with legs $h$ and $R - r$:
 \[l = \sqrt{(R - r)^2 + h^2}\]

 For an inscribed circle in an isosceles trapezoid, the slant side equals the sum of the two parallel radii drawn to the point of tangency. Hence:
 \[l = R + r\]

 Equating the two expressions for $l$:
 \[R + r = \sqrt{(R - r)^2 + (2a)^2}\]

 Squaring both sides:
 \[(R + r)^2 = (R - r)^2 + 4a^2\]
 \[R^2 + 2Rr + r^2 = R^2 - 2Rr + r^2 + 4a^2\]
 \[4Rr = 4a^2 \implies a = \sqrt{Rr}\]

 Thus,
 \[h = 2\sqrt{Rr}\]

 The volume of the frustum is:
 \[V_f = \frac{\pi h}{3}(R^2 + Rr + r^2) = \frac{2\pi\sqrt{Rr}}{3}(R^2 + Rr + r^2)\]

 The volume of the sphere is:
 \[V_s = \frac{4}{3}\pi a^3 = \frac{4}{3}\pi (Rr)^{3/2}\]

 The given condition $V_f = 2V_s$ becomes:
 \[\frac{2\pi\sqrt{Rr}}{3}(R^2 + Rr + r^2) = \frac{8\pi}{3}(Rr)^{3/2}\]

 Multiplying by $\frac{3}{\pi}$:
 \[2\sqrt{Rr}(R^2 + Rr + r^2) = 8(Rr)^{3/2}\]

 Dividing by $2\sqrt{Rr}$:
 \[R^2 + Rr + r^2 = 4Rr\]
 \[R^2 - 3Rr + r^2 = 0\]

 Dividing by $r^2$ and setting $x = \frac{R}{r}$:
 \[x^2 - 3x + 1 = 0\]
 \[x = \frac{3 \pm \sqrt{9 - 4}}{2} = \frac{3 \pm \sqrt{5}}{2}\]

 Since $R > r$ (the bottom base is larger), we take the larger root:
 \[\frac{R}{r} = \frac{3 + \sqrt{5}}{2}\]

 The final ratio is $\boxed{\textbf{(E) } \frac{3 + \sqrt{5}}{2}}$.

 (Note: This solution is still fundamentally similar to the previous ones)

 


	\item (\textit{AMC 10 10B 2014}) The numbers 1, 2, 3, 4, 5 are to be arranged in a circle. An arrangement is \textit{bad} if it is not true that for every $n$ from $1$ to $15$ one can find a subset of the numbers that appear consecutively on the circle that sum to $n$. Arrangements that differ only by a rotation or a reflection are considered the same. How many different bad arrangements are there?


 $\textbf {(A) } 1 \qquad \textbf {(B) } 2 \qquad \textbf {(C) } 3 \qquad \textbf {(D) } 4 \qquad \textbf {(E) } 5$


 
 \textbf{Answer: }B



\textbf{Solution 1}

We see that there are $5!$ total ways to arrange the numbers. However, we can always rotate these numbers so that, for example, the number $1$ is always at the top of the circle. Thus, there are only $4!$ ways under rotation. Every case has exactly $1$ reflection, so that gives us only $4!/2$, or $12$ cases, which is not difficult to list out. We systematically list out all $12$ cases. 

 Now, we must examine if they satisfy the conditions. We can see that by choosing one number at a time, we can always obtain subsets with sums $1, 2, 3, 4,$ and $5$. By choosing the full circle, we can obtain $15$. By choosing everything except for $1, 2, 3, 4,$ and $5$, we can obtain subsets with sums of $10, 11, 12, 13,$ and $14$. 

 This means that we now only need to check for $6, 7, 8,$ and $9$. However, once we have found a set summing to $6$, we can choose all remaining numbers and obtain a set summing to $15-6=9$, and similarly for $7$ and $8$. Thus, we only need to check each case for whether or not we can obtain $6$ or $7$. 

 We can make $6$ by having $4, 2$, or $3, 2, 1$, or $5, 1$. We can start with the group of three. To separate $3, 2, 1$ from each other, they must be grouped two together and one separate, like this.

 
                    \begin{figure}[H]
                        \centering
                        \includegraphics[width=0.5\textwidth]{Solution/AMC_10_10B_2014_24_Solution_Figure_1.png}
                    \end{figure}
                    

 Now, we note that $x$ is next to both blank spots, so we can't have a number from one of the pairs. So since we can't have $1$, because it is part of the $5, 1$ pair, and we can't have $2$ there, because it's part of the $4, 2$ pair, we must have $3$ inserted into the $x$ spot. We can insert $1$ and $2$ in $y$ and $z$ interchangeably, since reflections are considered the same.

 
                    \begin{figure}[H]
                        \centering
                        \includegraphics[width=0.5\textwidth]{Solution/AMC_10_10B_2014_24_Solution_Figure_2.png}
                    \end{figure}
                    

 We have $4$ and $5$ left to insert. We can't place the $4$ next to the $2$ or the $5$ next to the $1$, so we must place $4$ next to the $1$ and $5$ next to the $2$.

 
                    \begin{figure}[H]
                        \centering
                        \includegraphics[width=0.5\textwidth]{Solution/AMC_10_10B_2014_24_Solution_Figure_3.png}
                    \end{figure}
                    

 This is the only solution to make $6$ "bad."

 Next we move on to $7$, which can be made by $3, 4$, or $5, 2$, or $4, 2, 1$. We do this the same way as before. We start with the three group. Since we can't have $4$ or $2$ in the top slot, we must have one there, and $4$ and $2$ are next to each other on the bottom. When we have $3$ and $5$ left to insert, we place them such that we don't have the two pairs adjacent.

 
                    \begin{figure}[H]
                        \centering
                        \includegraphics[width=0.5\textwidth]{Solution/AMC_10_10B_2014_24_Solution_Figure_4.png}
                    \end{figure}
                    

 This is the only solution to make $7$ "bad."

 We've covered all needed cases, and the two examples we found are distinct, therefore the answer is $\boxed{\textbf {(B) }2}$.

 



\textbf{Solution 2 (Complex Compliments)}

We can use a clever approach for this problem by finding out how many non-adjacency's exist and subtracting it by the total configurations.

 The formula \( C_n(A) = \frac{n}{k} \binom{n-k-1}{k-1} \) represents the Circular Gap Theorem, a crucial formula for finding out how many compliment cases exist where we want \( k \) values to not be adjacent from \( n \) total values.

 Considering reflections and rotations are practically the same, we see that we have \( (n-1)!/2 \), \( n=5 \) ways = \( 4!/2 = 24/2 = 12 \) total configurations.

 Now, we do some casework. Notice that we have \( n=5 \) values. Our \( k \in \{1,2,3,4,5\} \). We see that the values \( k \in \{3, 4, 5\} \) don't work as they make \( \binom{n-k-1}{k-1} \) invalid (3 gives \( 5-3-1 \choose 3-1 \) = \( 1 \choose 2 \) which doesn't exist, 4 gives \( 0 \choose 3 \) which doesn't exist, and 5 gives \( -1 \choose 4 \) which obviously doesn't exist).

 We solve for \( k \in \{1,2\} \). When \( k \in 1 \), we get \( C_n(A) = \frac{5}{1} \binom{5-1-1}{1-1} = 5(1) = 5 \) cases and when \( k \in 2 \), we get \( C_n(A) = \frac{5}{2} \binom{5-2-1}{2-1} = \frac{5}{2}(2) = 5 \). The total cases that lead to a good sum is \( 5+5=10 \).

 The answer isn't 10 however, because the theorem finds the compliments of the values, so the number of compliments is 10. We need to find the bad cases, which is \( 12 - 10 \) or $\textbf{(B)} \text{ }2$

 \textbf{Circular Gap Theorem Simplification} It is also helpful to note \( C_n(A) = \frac{n}{k} \binom{n-k-1}{k-1} \) = \( \binom{n}{2k} \).

 



\textbf{Solution 3}

Note that any ordering satisfies the following numbers:

 $1$ through $5:$ choose the number

 $10$ through $14:$ choose all numbers excluding a specific one (such as $1, 2, 3, 4$ in some order for $10$)

 $15:$ choose all the numbers.

 Now, the pairs $7, 8$ and $6, 9$ are the same cases, since if a sequence satisfies a number, we can choose all the remaining numbers to make the other number. ($1, 2, 3$ for $6$, then $4, 5$ for $9$.)

 Thus, we have two cases, whether a sequence doesn't make $6$ or whether a sequence doesn't make $7.$

 $6$ can be made by $(1,2,3), (1, 5), (2, 4).$ We can put $1, 2, 3$ around the circle. $4$ and $5$ now need to go in $2$ of the $3$ spots in between $1, 2, 3.$ Also keeping in mind the other two ways to make $6,$ $4$ has to go in the spot opposite of $2$ and $5$ has to go in the spot opposite of $1.$ Thus the only ordering that works is $1, 4, 3, 5, 2$ (ignore rotations and reflections).

 Similarly, for the case with $7,$ the only ordering that works is $1, 5, 4, 2, 3$ with gives the answer of $\mathbf{(B) }2.$

 



\textbf{Solution 4 (Minimal Casework)}

Note that $\{1,2,3,4,5,10,11,12,13,14,15\}$ will always be there. Thus, we need to prevent either the $(6,9)$ or the $(7,8)$ pair. We consider the neighbors of 5. 

 Case 1: No 7s. Then we have $1-5-4$, and the neighbor of 4 cannot be 3 so the full config must be $3-1-5-4-2$. 

 Case 2: No 6s. Then we have $2-5-3$, and the neighbor of 2 cannot be 4, so the full config must be $1-2-5-3-4$. 

 Both of these are bad, and they're the only bads, so the answer is $\textbf{(B)} \text{ }2$.

 



\textbf{Solution 5 (Maximal Casework)}

Like all previous solutions, we only need to find an arrangement where a sum of $6$ or $7$ isn't achievable. Reflections and rotations are considered the same, so there are actually only $\frac{5!}{5} \cdot \frac{1}{2} = 12$ distinct arrangements. We can simply draw all $12$ arrangements and cross out the ones that aren't bad, leaving us with $\boxed{2}$ bad arrangements.

 


	\item (\textit{AMC 10 10B 2014}) In a small pond there are eleven lily pads in a row labeled $0$ through $10$. A frog is sitting on pad $1$. When the frog is on pad $N$, $0<N<10$, it will jump to pad $N-1$ with probability $\frac{N}{10}$ and to pad $N+1$ with probability $1-\frac{N}{10}$. Each jump is independent of the previous jumps. If the frog reaches pad $0$ it will be eaten by a patiently waiting snake. If the frog reaches pad $10$ it will exit the pond, never to return. What is the probability that the frog will escape without being eaten by the snake?


 $\textbf {(A) } \frac{32}{79} \qquad \textbf {(B) } \frac{161}{384} \qquad \textbf {(C) } \frac{63}{146} \qquad \textbf {(D) } \frac{7}{16} \qquad \textbf {(E) } \frac{1}{2}$


 
 \textbf{Answer: }C



\textbf{Solution 1}

A long, but straightforward bash:

 Define $P(N)$ to be the probability that the frog survives starting from pad N.

 Then note that by symmetry, $P(5) = 1/2$, since the probabilities of the frog moving subsequently in either direction from pad 5 are equal.

 We therefore seek to rewrite $P(1)$ in terms of $P(5)$, using the fact that

 $P(N) = \frac {N} {10}P(N - 1) + \frac {10 - N} {10}P(N + 1)$

 as said in the problem.

 Hence $P(1) = \frac {1} {10}P(0) + \frac {9} {10}P(2) = \frac {9} {10}P(2)$

 $\Rightarrow P(2) = \frac {10} {9}P(1)$

 Returning to our original equation:

 $P(1) = \frac {9} {10}P(2) = \frac {9} {10}\left(\frac{2} {10}P(1) + \frac{8} {10}P(3)\right)$

 $= \frac {9} {50}P(1) + \frac {18} {25}P(3) \Rightarrow P(1) - \frac {9} {50}P(1)$ $= \frac {18} {25}P(3)$

 $\Rightarrow P(3) = \frac {41} {36}P(1)$

 Returning to our original equation:

 $P(1) = \frac {9} {50}P(1) + \frac {18} {25}\left(\frac {3} {10}P(2) + \frac {7} {10}P(4)\right)$

 $= \frac {9} {50}P(1) + \frac {27} {125}P(2) + \frac {63} {125}P(4)$

 $= \frac {9} {50}P(1) + \frac {27} {125}\left(\frac {10} {9}P(1)\right) + \frac {63} {125}\left(\frac {4} {10}P(3) + \frac {6} {10}P(5)\right)$

 Cleaning up the coefficients, we have: 

 $= \frac {21} {50}P(1) + \frac {126} {625}P(3) + \frac {189} {625}P(5)$

 $= \frac {21} {50}P(1) + \frac {126} {625}\left(\frac {41} {36}P(1)\right) + \frac {189} {625}\left(\frac {1} {2}\right)$

 Hence, $P(1) = \frac {525} {1250}P(1) + \frac {287} {1250}P(1) + \frac {189} {1250}$

 $\Rightarrow P(1) - \frac {812} {1250}P(1) = \frac {189} {1250} \Rightarrow P(1) = \frac {189} {438}$

 $= \boxed{\frac {63} {146}\, (C)}$

 Or set $P(1)=a,P(2)=b,P(3)=c,P(4)=d,P(5)=e=1/2$:
 \[a=0.1\emptyset+0.9b,b=0.2a+0.8c,c=0.3b+0.7d,d=0.4c+0.6e\]
 \[10a=\emptyset+9b,10b=2a+8c,10c=3b+7d,10d=4c+6e\]
 \[\implies b=\frac{10a-\emptyset}{9},c=\frac{5b-a}{4},d=\frac{10c-3b}{7},e=\frac{5d-2c}{3}=1/2\]$b=\frac{10a}{9}$

 $c=\frac{5\left(\frac{10a}{9}\right)-a}{4}=\frac{\frac{50a}{9}-a}{9}=\frac{41a}{36}$

 $d=\frac{10\left(\frac{41a}{36}\right)-3\left(\frac{30a}{9}\right)}{7}=\frac{\frac{205a}{18}-\frac{10a}{3}}{7}=\frac{145a}{126}$

 $e=\frac{5\left(\frac{145a}{126}\right)-2\left(\frac{41a}{36}\right)}{3}=\frac{\frac{725a}{126}-\frac{41a}{18}}{3}=\frac{73a}{63}$

 Since $e=\frac{1}{2}$, $\frac{73a}{63}=\frac{1}{2}\implies a=\boxed{\textbf{(C) }\frac{63}{146}}$.

 



\textbf{Solution 2}

Notice that the probabilities are symmetrical around the fifth lily pad. If the frog is on the fifth lily pad, there is a $\frac{1}{2}$ chance that it escapes and a $\frac{1}{2}$ chance that it gets eaten. Now, let $P_k$ represent the probability that the frog escapes if it is currently on pad $k$. We get the following system of $5$ equations:
 \[P_1=\frac{9}{10}\cdot P_2\]
 \[P_2=\frac{2}{10}\cdot P_1 + \frac{8}{10}\cdot P_3\]
 \[P_3=\frac{3}{10}\cdot P_2 + \frac{7}{10}\cdot P_4\]
 \[P_4=\frac{4}{10}\cdot P_3 + \frac{6}{10}\cdot P_5\]
 \[P_5=\frac{5}{10}\]We want to find $P_1$, since the frog starts at pad $1$. Solving the above system yields $P_1=\frac{63}{146}$, so the answer is $\boxed{(C)}$.

 System of equations (alligator112): Start by removing all fractions. From here on out, we will use substitution while avoiding fractions.
 \[10P_1 = 9P_2\]
 \[10P_2 = 2P_1 + 8P_3\]
 \[10P_3 = 3P_2 + 7P_4\]
 \[10P_4 = 4P_3 + 6P_5 = 4P_3 + 3\]
 \[100 P_3 = 30P_2 + 70P_4 = 30P_2 + 7(4P_3 + 3) \rightarrow 72P_3 = 30P_2 + 21\]
 \[90P_2 = 18P_1 + 72P_3 = 18P_1 + 30P_2 + 21 \rightarrow 60P_2 = 18P_1 + 21\]
 \[200P_1 = 180P_2 = 3(18P_1 + 21) = 54P_1 + 63 \rightarrow P_1 = \frac{63}{146}\]

 



\textbf{Solution 3}

Assign each lily pad a value, with pad $0$ having value $0$ and pad $i > 0$ having value $\sum_{k=0}^{i-1}\frac{1}{\binom{9}{k}}$. If we treat the process as a random walk $X_0, X_1, \cdots$ over these values, note that $\mathbb{E}[X_{n+1}] = n$ for all $n \ge 0$. This is purely by construction, as the expected gain at pad $k$ is $\binom{9}{k}\cdot \frac{10-k}{10} - \binom{9}{k-1}\cdot \frac{k}{10} = 0$. Therefore the process is a martingale. Since the stopping time has finite expectation and the increments are bounded, the \href{https://en.wikipedia.org/wiki/Optional\_stopping\_theorem}{Optional Stopping Thoerem} applies and the expected final outcome is $0$. This gives the probability of reaching safety as $\left(\sum_{k=0}^{9}\frac{1}{\binom{9}{k}}\right)^{-1} = \frac{63}{146}$, which is choice $\boxed{(C)}$.

 


 



\textbf{Solution 4 (Explaining the Martingale Shortcut Mentioned Previously)}

Instead of grinding through a massive system of 9 equations, we can solve this problem by assigning a "score value" to each lily pad. 

 Normally, if a game is perfectly fair (like a 50/50 coin flip), your average score after a turn is the same as your score before the turn. But this frog pond is rigged because the probabilities change at every pad. To fix this, we can mathematically design a custom scoring system that forces the game to be perfectly fair. 

 Once the scoring system is perfectly balanced, the law of expected value tells us that the probability of winning is just a simple ratio:
 \[\text{Escape Probability} = \frac{\text{Score of your Starting Pad}}{\text{Score of the Winning Pad (Pad 10)}}\]

 
                \begin{itemize}
                    \item (Note: We set the score of Pad 0 to be 0 points, so getting eaten by the snake gives 0 credit).*
                \end{itemize}


                \textbf{Step 1: Find the Step Rule} To keep our scores perfectly balanced at any pad $N$, the reward for a rare forward jump must perfectly offset the risk of an easy backward jump. This gives us a simple rule for how the point-gaps between pads must scale:

 
 \[\text{Step Multiplier} = \frac{\text{Probability of jumping Left}}{\text{Probability of jumping Right}} = \frac{\frac{N}{10}}{\frac{10-N}{10}} = \frac{N}{10-N}\]

 This fraction $\frac{N}{10-N}$ tells us exactly how much larger the next point-step must be to keep our scoring scale balanced.

 \textbf{Step 2: Build the Chain of Gaps} We start by declaring that moving from Pad 0 to Pad 1 awards exactly $1$ point. To find the point value of each subsequent step, we take our current step value and multiply it by our Step Multiplier rule:

 
                \begin{itemize}
                    \item \textbf{Step 1 (Pad 0 to 1):} Worth $1$ point. Pad 1 Score = $1$.\item \textbf{Step 2 (Pad 1 to 2):} Multiply by the rule at $N=1 \implies 1 \times \frac{1}{9} = \frac{1}{9}$ points.\item \textbf{Step 3 (Pad 2 to 3):} Multiply by the rule at $N=2 \implies \frac{1}{9} \times \frac{2}{8} = \frac{1}{36}$ points.\item \textbf{Step 4 (Pad 3 to 4):} Multiply by the rule at $N=3 \implies \frac{1}{36} \times \frac{3}{7} = \frac{1}{84}$ points.\item \textbf{Step 5 (Pad 4 to 5):} Multiply by the rule at $N=4 \implies \frac{1}{84} \times \frac{4}{6} = \frac{1}{126}$ points.
                \end{itemize}


                Because the pond's layout is perfectly symmetric, the point values for the remaining five steps mirror these numbers exactly in reverse order. This gives us all 10 step values:
 \[\left[ 1, \, \frac{1}{9}, \, \frac{1}{36}, \, \frac{1}{84}, \, \frac{1}{126}, \, \frac{1}{126}, \, \frac{1}{84}, \, \frac{1}{36}, \, \frac{1}{9}, \, 1 \right]\]

 \textbf{Step 3: Find the Total Score and Solve} Now we can easily calculate our two target scores:

 
            \begin{enumerate}
                \item \textbf{Score of our Starting Pad (Pad 1):} This is just the value of the very first step: $1$.\item \textbf{Score of the Winning Pad (Pad 10):} This is the total sum of all 10 steps combined. Grouping the identical matching pairs over a common denominator of 126 makes this simple to calculate:
            \end{enumerate}
            
 \[\text{Total Score of Pad 10} = 2 \left(1 + \frac{1}{9} + \frac{1}{36} + \frac{1}{84} + \frac{1}{126}\right) = 2 \left(\frac{252 + 28 + 7 + 3 + 2}{126}\right) = \frac{146}{63}\]

 Finally, we divide our starting score by the winning pad's score to get the final probability:

 
 \[P = \frac{\text{Starting Score}}{\text{Winning Score}} = \frac{1}{\frac{146}{63}} = \boxed{\textbf{(C) } \frac{63}{146}}\]

 \textbf{Remark (Markov Chain)} 
            \begin{enumerate}[label=(\roman*)]
                \item <i>Main article: \href{https://artofproblemsolving.com/wiki/index.php?title=Markov\_Chains}{Markov Chains}</i>
            \end{enumerate}
            We can represent this problem with the following State Transition Diagram of \href{https://en.wikipedia.org/wiki/Absorbing\_Markov\_chain}{Absorbing Markov Chain}. State $0$ and state $10$ are the absorbing states. This problem asks for the \href{https://en.wikipedia.org/wiki/Absorbing\_Markov\_chain\#Absorbing\_probabilities}{Absorbing Probability} of state $10$.

 
                    \begin{figure}[H]
                        \centering
                        \includegraphics[width=0.5\textwidth]{Solution/AMC_10_10B_2014_25_Solution_Figure_1.png}
                    \end{figure}
                    Let $p_{ij} = P(X_{n+1} = j | X_n = i)$, the probability that state $i$ transits to state $j$ on the next step. 

 Let $a_i$ be the absorbing probability of being absorbed in the absorbing state when starting from transient state $i$.

 $a_i = \sum_{j} (p_{ij} \cdot a_{j})$

 $a_i$ is the sum of the products of $p_{ij}$ and $a_j$ of all the next state $j$.

 Note:

 This sum results in the same system of equations as solution 2, and the solution follows as well. 

 


\end{enumerate}